\documentclass[preview]{standalone}
\usepackage{amsmath}
\usepackage{amssymb}
\usepackage{stellar}
\usepackage{definitions}
\usepackage{bettelini}
\begin{document}
\id{diffeq-linear-systems}
\genpage
\section{Linear systems}

\begin{snippetdefinition}{diffeq-linear-system-definition}{Linear differential system}
    Let \(I\) be an open interval,
    \(A\in C(I,\realnumbers^{m\times m})\), \(b\in C(I,\realnumbers^m)\).
    A \emph{linear differential system} is \(y'=A(t)y+b(t)\).
    Its associated homogeneous system is \(y'=A(t)y\).
    The linear operator
    \[
        \Lambda\colon C^1(I,\realnumbers^m)\fromto C(I,\realnumbers^m),
        \qquad \Lambda u=u'-Au
    \]
    has the homogeneous solution space as its kernel.
\end{snippetdefinition}

\begin{snippetcorollary}{diffeq-linear-system-global}{Global solvability of linear systems}
    If \(A\in C(I,\realnumbers^{m\times m})\) and \(b\in C(I,\realnumbers^m)\)
    on an open interval, every problem \(y'=Ay+b\), \(y(t_0)=y_0\),
    has a unique solution on all of \(I\).
\end{snippetcorollary}

\begin{snippetproof}{diffeq-linear-system-global-proof}{diffeq-linear-system-global}{Global solvability of linear systems}
    The operator norm gives
    \[
        \|A(t)u-A(t)v\|\leq\|A(t)\|\|u-v\|,\qquad
        \|A(t)y+b(t)\|\leq\|b(t)\|+\|A(t)\|\|y\|.
    \]
    Continuity of the coefficients gives local uniform Lipschitz bounds
    and a linear-growth bound. Apply the global growth theorem.
\end{snippetproof}

\begin{snippetproposition}{diffeq-linear-system-affine}{Affine solution space}
    Let \(A,b\) be continuous on an open interval \(I\), and let
    \(\widetilde y\) be one solution of \(y'=Ay+b\).
    The complete solution \set is \(\widetilde y+\ker\Lambda\),
    where \(\Lambda u=u'-Au\).
\end{snippetproposition}

\begin{snippetproof}{diffeq-linear-system-affine-proof}{diffeq-linear-system-affine}{Affine solution space}
    If \(\Lambda y=\Lambda\widetilde y=b\), linearity gives
    \(\Lambda(y-\widetilde y)=0\).
    Conversely \(\Lambda(\widetilde y+w)=b\) for every \(w\in\ker\Lambda\).
\end{snippetproof}

\begin{snippettheorem}{diffeq-linear-system-dimension}{Dimension of the homogeneous solution space}
    For \(A\in C(I,\realnumbers^{m\times m})\), the solutions of \(y'=Ay\)
    on the open interval \(I\) form a vector space of dimension \(m\).
    Evaluation \(y\mapsto y(t_0)\) at any \(t_0\in I\) is an isomorphism
    to \(\realnumbers^m\).
\end{snippettheorem}

\begin{snippetproof}{diffeq-linear-system-dimension-proof}{diffeq-linear-system-dimension}{Dimension of the homogeneous solution space}
    Choose a basis \(e_1,\ldots,e_m\) of \(\realnumbers^m\).
    Let \(\varphi_i\) solve \(\varphi_i'=A\varphi_i\), \(\varphi_i(t_0)=e_i\).
    If \(\sum c_i\varphi_i=0\), evaluation at \(t_0\) gives
    \(\sum c_ie_i=0\), hence all \(c_i=0\).
    Conversely, for any homogeneous solution \(\varphi\), expand
    \(\varphi(t_0)=\sum c_ie_i\). The solution \(\psi=\sum c_i\varphi_i\)
    has the same initial value; uniqueness gives \(\psi=\varphi\).
    Thus these \(m\) solutions are an independent spanning family.
\end{snippetproof}

\begin{snippetdefinition}{diffeq-fundamental-matrix-definition}{Fundamental matrix}
    For \(y'=A(t)y\), with continuous \(A\) on an open interval \(I\),
    a \emph{fundamental matrix} is \(\Phi=[\varphi_1,\ldots,\varphi_m]\),
    whose columns form a basis of homogeneous solutions.
    It satisfies \(\Phi'=A\Phi\); its determinant \(W=\det\Phi\) is the
    \emph{Wronskian}. It is \emph{normalized at \(t_0\)} if \(\Phi(t_0)=I_m\).
\end{snippetdefinition}

\begin{snippettheorem}{diffeq-system-wronskian-independence}{Independence and the Wronskian}
    Let \(\varphi_1,\ldots,\varphi_m\) solve \(y'=A(t)y\) on an open interval,
    with \(A\) continuous. The following are equivalent:
    they are linearly independent as \function[functions];
    \(\det[\varphi_1(t),\ldots,\varphi_m(t)]\neq0\) at one point;
    this determinant is nonzero at every point.
\end{snippettheorem}

\begin{snippetproof}{diffeq-system-wronskian-independence-proof}{diffeq-system-wronskian-independence}{Independence and the Wronskian}
    A constant nonzero vector \(c\) with \(\sum c_i\varphi_i=0\) gives a
    nontrivial null vector for the column matrix at every point.
    Conversely, if its determinant vanishes at \(t_*\), fix a nonzero
    \(c\) in that matrix's kernel. The \function \(w=\sum c_i\varphi_i\)
    solves \(w'=Aw\) and \(w(t_*)=0\); uniqueness forces \(w\equiv0\).
    It is essential that these coefficients are fixed at \(t_*\), not
    chosen afresh as functions of time.
\end{snippetproof}

\begin{snippetexample}{diffeq-generic-zero-determinant}{Why the common equation matters}
    The vector-valued \function[functions]
    \(\varphi_1(x)=(x^2,2x)\), \(\varphi_2(x)=(x|x|,2|x|)\)
    are independent on \(\realnumbers\), since they agree on the positive
    half-line and are negatives on the negative half-line.
    Their column determinant nevertheless vanishes everywhere.
    They are not solutions of one common continuous linear system.
    For arbitrary \function[functions], dependence persists on restriction,
    and independence persists on extension, but neither converse holds.
    For solutions of a common system, the Wronskian criterion gives
    independence on every nonempty open subinterval.
\end{snippetexample}

\begin{snippettheorem}{diffeq-abel-liouville-identity}{Abel--Liouville identity}
    Let \(A\in C(I,\realnumbers^{m\times m})\), and let the columns of
    \(\Phi\in C^1(I,\realnumbers^{m\times m})\) solve \(y'=Ay\), without
    assuming their independence. Then
    \[
        W'=(\operatorname{tr}A)W,\qquad
        W(t)=W(t_0)\exp\left(\int_{t_0}^t\operatorname{tr}A(s)\,ds\right),
        \quad W=\det\Phi.
    \]
    In particular \(W\) is either identically zero or never zero.
\end{snippettheorem}

\begin{snippetproof}{diffeq-abel-liouville-identity-proof}{diffeq-abel-liouville-identity}{Abel--Liouville identity}
    Differentiate the determinant by multilinearity in its rows.
    The \(i\)-th differentiated row is
    \(\sum_j A_{ij}\) times the \(j\)-th row of \(\Phi\).
    Each term with \(j\neq i\) has repeated rows and vanishes; the remaining
    term is \(A_{ii}\det\Phi\). Summing yields \(W'=(\operatorname{tr}A)W\).
    The scalar integrating-factor formula gives the exponential expression.
\end{snippetproof}

\section{Variation of constants}

\begin{snippettheorem}{diffeq-system-variation-constants}{Variation of constants for systems}
    Let \(A,b\) be continuous on an open interval \(I\), and \(\Phi\) be a
    fundamental matrix of \(y'=Ay\). Every solution of \(y'=Ay+b\) is
    \[
        y(t)=\Phi(t)\left(c+\int_{t_0}^t\Phi(s)^{-1}b(s)\,ds\right).
    \]
    For the initial condition \(y(t_0)=y_0\), choose \(c=\Phi(t_0)^{-1}y_0\).
\end{snippettheorem}

\begin{snippetproof}{diffeq-system-variation-constants-proof}{diffeq-system-variation-constants}{Variation of constants for systems}
    Invertibility makes every \(C^1\) \function expressible as
    \(y=\Phi c(t)\), with \(c(t)=\Phi(t)^{-1}y(t)\in C^1\).
    Substitute and use \(\Phi'=A\Phi\):
    \[
        \Phi'c+\Phi c'=A\Phi c+b,\qquad
        \Phi c'=b,\qquad c'=\Phi^{-1}b.
    \]
    Integration gives the formula. Thus a particular solution is
    \(\Phi(t)\int_{t_0}^t\Phi^{-1}b\); the homogeneous solutions are
    \(\Phi(t)c\) with constant \(c\), and the complete solution space is affine of dimension \(m\).
    For \(m=1\), \(\Phi=e^{\int_{t_0}^t a}\) gives
    \[
        y(t)=e^{\int_{t_0}^t a(s)\,ds}
        \left(c+\int_{t_0}^t b(s)e^{-\int_{t_0}^s a(r)\,dr}\,ds\right),
    \]
    recovering the scalar formula.
\end{snippetproof}

\section{Constant matrices and scalar equations}

\begin{snippetdefinition}{diffeq-matrix-exponential-definition}{Matrix exponential}
    For a square matrix \(A\), its \emph{exponential} is
    \(e^A=\sum_{j=0}^{\infty}A^j/j!\).
\end{snippetdefinition}

\begin{snippetproposition}{diffeq-constant-system-exponential}{Constant-coefficient fundamental matrix}
    For a constant real square matrix \(A\), the exponential series converges
    absolutely, and \(\Phi(t)=e^{(t-t_0)A}\) is the normalized fundamental
    matrix of \(y'=Ay\).
\end{snippetproposition}

\begin{snippetproof}{diffeq-constant-system-exponential-proof}{diffeq-constant-system-exponential}{Constant-coefficient fundamental matrix}
    A submultiplicative matrix norm gives
    \(\sum\|A^j\|/j!\leq\sum\|A\|^j/j!<\infty\).
    The series and its derivative converge uniformly on compact time intervals,
    so \(\Phi'=A\Phi\), \(\Phi(t_0)=I_m\).
    Its columns form a basis by their initial values.
    In general, for time-varying \(A(t)\), simply exponentiating \(\int A\)
    does not give a fundamental matrix; noncommutativity matters.
\end{snippetproof}

\begin{snippetproposition}{diffeq-companion-system}{Scalar linear equations as systems}
    For continuous \(a_1,\ldots,a_m,b\), the equation
    \[
        y^{(m)}=a_1y+a_2y'+\cdots+a_my^{(m-1)}+b
    \]
    becomes \(Y'=AY+B\), with \(Y=(y,\ldots,y^{(m-1)})\), \(B=(0,\ldots,0,b)\), and
    \[
        A=\begin{pmatrix}
        0&1&0&\cdots&0\\
        0&0&1&\cdots&0\\
        \vdots&&&\ddots&\vdots\\
        0&0&0&\cdots&1\\
        a_1&a_2&a_3&\cdots&a_m
        \end{pmatrix}.
    \]
    For \(m=1\), take \(A=(a_1)\). Thus all initial problems are globally
    solvable on the coefficient interval, and the homogeneous space has dimension \(m\).
    The columns \((w_j,w_j',\ldots,w_j^{(m-1)})\) form precisely the scalar
    Wronskian matrix. When the coefficients are constant, the ODE characteristic
    polynomial is the characteristic polynomial of this companion matrix.
\end{snippetproposition}

\begin{snippetexample}{diffeq-exponential-vandermonde}{Independence of exponential solutions}
    For distinct \(\alpha_1,\ldots,\alpha_m\), the Wronskian of
    \(e^{\alpha_1x},\ldots,e^{\alpha_mx}\) is
    \[
        e^{(\alpha_1+\cdots+\alpha_m)x}
        \det(\alpha_j^{i-1})_{i,j=1}^m
        =e^{(\alpha_1+\cdots+\alpha_m)x}
           \prod_{i<j}(\alpha_j-\alpha_i)\neq0.
    \]
    For the repeated pair \(e^{\alpha x},xe^{\alpha x}\), it is
    \(e^{2\alpha x}\neq0\). Polynomial factors with a common exponential
    are also independent by their different degrees.
    These analytic \function[functions] remain independent on every open
    subinterval: a vanishing linear combination there vanishes everywhere.
    For complex roots one uses real and imaginary parts; growth order alone
    does not distinguish the sine and cosine solutions.
\end{snippetexample}
\end{document}
